\paragraph{Purpose}
A feed-forward network (or \emph{multi-layer perceptron}, MLP) defines a mapping
$\bm{y} = f(\bm{x};\bm{\theta})$ and learns the parameters $\bm{\theta}$ that give the best
approximation of an unknown function $f^{*}$. \tB{Information flows from the input to the
output with no feedback connection}, otherwise the network would be recurrent.

\paragraph{Theory}
The network is a composition of $l$ layers, $f = f_{l}\circ\cdots\circ f_{1}$, each one
being an affine map followed by a non-linear activation function $g$ applied element-wise:
\begin{center}
    \frB{$\bm{h}_{k} = g\left(\bm{W}_{k}\bm{h}_{k-1} + \bm{b}_{k}\right),\quad
    \bm{h}_{0} = \bm{x},\quad \hat{\bm{y}} = \bm{h}_{l}$}
\end{center}
Without the non-linearity, the composition of affine maps collapses into a \uB{single
affine map}, so depth would bring nothing.

\paragraph{Universal approximation theorem}
A feed-forward network with \uB{one hidden layer}, a non-polynomial activation function and
enough hidden units can approximate any continuous function on a compact set
(Cybenko 1989, Hornik 1991). The theorem says nothing about \tB{how many units are
needed}, nor about whether gradient descent can \tB{find} the right parameters.

\paragraph{Activation functions}
\begin{center}
\begin{tabular}{|l|c|l|}
    \hline
    Name & $g(x)$ & Remark\\
    \hline
    Sigmoid & $\sigma(x) = \dfrac{1}{1+e^{-x}}$ & saturates, $\sigma'\leq\frac{1}{4}$: gradients vanish\\
    Tanh & $\tanh(x)$ & zero-centered, but also saturates\\
    ReLU & $\max(0, x)$ & gradient $0$ or $1$, cheap, units can ``die''\\
    GELU & $x\,\Phi(x)$ & smooth ReLU, used in GPT and BERT\\
    \hline
\end{tabular}
\end{center}
($\Phi$ is the cdf of the standard normal distribution.) Saturating activations multiply
gradients by factors smaller than one at each layer in the backward pass, which is the same
phenomenon as in the long-term dependencies of recurrent networks.

\paragraph{Output layer and loss}
The last layer and the loss are chosen from the probabilistic model of the target,
the loss being the negative log-likelihood:
\begin{center}
\begin{tabular}{|l|l|l|}
    \hline
    Task & Output layer & Loss\\
    \hline
    Regression & linear & squared error (Gaussian likelihood)\\
    Binary classification & sigmoid & binary cross-entropy (Bernoulli)\\
    Multi-class classification & softmax & cross-entropy (categorical)\\
    \hline
\end{tabular}
\end{center}

\paragraph{Weaknesses}
\begin{itemize}
    \item no parameter sharing: the number of parameters grows with the input size;
    \item fixed input size and no notion of order or locality, which motivates
        convolutional networks (grids) and recurrent networks or Transformers (sequences).
\end{itemize}

\paragraph{Relationships with other methods}
\begin{itemize}
    \item with no hidden layer, a sigmoid (resp. softmax) output gives logistic regression;
    \item the position-wise feed-forward sub-layer of the Transformer is a two-layer MLP
        applied independently to each token.
\end{itemize}
